\subsection{由边界条件认识能级离散}
以一维无限深势阱 $0<x<a$ 为模型：阱内 $V=0$，阱外概率为零，波函数在两壁满足 $\psi(0)=\psi(a)=0$。定态方程化为
\[
 \psi''+k^2\psi=0,\qquad E=\frac{\hbar^2k^2}{2m}.
\]
第一边界排除余弦项，第二边界要求 $\sin(ka)=0$。非零波函数对应 $k_n=n\pi/a$，$n=1,2,\ldots$，归一化后
\[
 \psi_n(x)=\sqrt{\frac2a}\sin\frac{n\pi x}{a},\qquad
 E_n=\frac{n^2\pi^2\hbar^2}{2ma^2}=\frac{n^2h^2}{8ma^2}.
\]
$n=0$ 只给处处为零的函数，不能归一化，所以基态能量不为零。这个模型说明离散值来自波动方程与边界条件，而不是任意把经典能量挑几个数值留下。
对电子和 $a=1.00\,\mathrm{nm}$，$E_1\approx0.376\,\mathrm{eV}$，$E_2-E_1=3E_1\approx1.13\,\mathrm{eV}$。阱宽减半时所有能级增为四倍。
无限势垒是理想边界；有限势垒会出现穿透尾部，不能继续强令阱外波函数处处为零。

衍射计算则先由加速能量求 $p$ 和 $\lambda=h/p$，再使用晶面的相位条件。加速电压、波长、散射角属于同一计算链；晶面夹角与法线夹角互为余角，使用布拉格式前应明确测角方式。
